\section{Conformance}
\label{sec:conformance}

Each result below is testable, and most are already tested. The table maps
every theorem to the specification section it formalises, the conformance
vectors of \SPEC{} Appendix B that exercise it, and the test in the
reference implementation
(\code{services/dispat/internal/plan/}) that stands behind it. The point of the
mapping is that a result with no vector and no test is a claim resting on the
proof alone, and the last column is where that shows.

\begin{center}
\footnotesize
\begin{longtable}{@{}L{3.0cm}L{1.5cm}L{1.8cm}L{8.0cm}@{}}
\toprule
Result & \SPEC{} & Vectors & Test \\
\midrule
\endhead
\cref{thm:g1} (\G{1}) & 13.7c & 67, 68 & \code{TestPropagationDepthIsShortestPath} \\
\cref{thm:g2} (\G{2}) & 13.7c & 66, 74, 76 & \code{TestCatchUpAfterConsumerFailure}, \code{TestCatchUpForNeverReleasedConsumer} \\
\cref{thm:g3} (\G{3}) & 13.7c & 70, 84, 137 & \code{TestCatchUpOnATrainWithPublishedOwnWork} \\
\cref{thm:g4} (\G{4}) & 13.7c & 66, 83, 133 & \code{TestCatchUpDischargesExactlyOnce} \\
\cref{thm:g5} (\G{5}) & 13.7c & 67, 68, 134 & \code{TestCatchUpNeverWidensBlastRadius} \\
\cref{thm:g6} (\G{6}) & 13.7c, 19.6 & 77, 77a, 133 & \code{TestPrereleaseTrainConvergesAfterRelease}, \code{TestPropagatedTrainWorkConverges} \\
\cref{prop:tagall} & 13.10a & 77, 77a & \code{TestFixedConvergesWhenNothingPending} \\
\cref{thm:g7} (\G{7}) & 13.7c, 9.3 & 60r, 97c & \code{TestChannelAxisConverges}, \code{TestTransitionIsIdempotent} \\
\cref{thm:g8} (\G{8}) & 13.7c, 9.2 & 94, 95, 97, 136 & \code{TestPropagatedStableNeverGraduates}, \code{TestCaretAloneDoesNotDragConsumersOntoATrain} \\
\cref{thm:determinism} & 17.2 & 82 & \code{TestConflictingDirectChannelsPickTheNewest}, \code{TestConflictingPropagatedChannelsPickTheNewest} \\
\cref{cor:idempotency} & 17.2 & 82, 83 & \code{TestNoCatchUpWhenConsumerIsFresh} \\
\cref{thm:resumption} & 17.2, 19.3 & 70, 84 & \code{TestCatchUpAfterConsumerFailure} \\
\cref{thm:duality} & 13.7b & 135 & \code{TestStaleSourcesAgreesWithPropagation} \\
\cref{prop:screen} & 13.7b & --- & --- \\
\cref{prop:order} & 19.2 & 69, 80a & --- (executor, outside \code{plan}) \\
\cref{prop:windows} & 13.11 & --- & \code{TestPlanningQueriesGitOncePerPackage} \\
\cref{prop:bucket} & 13.11 & 82c--82j, 138 & --- (no traversal reuse implemented) \\
\cref{prop:hoist} & 9.3a, 13.11 & 82c, 82d & \code{TestCaretAloneDoesNotDragConsumersOntoATrain} \\
\cref{prop:invert} & 13.8, 13.11 & 82e, 82f & \code{TestDirectChannelBeatsPropagated} \\
\cref{prop:parse} & 20.7 & 99 & \code{FuzzParse} (\code{pkg/ccme}) \\
\bottomrule
\end{longtable}
\end{center}

Three rows deserve comment.

\Cref{thm:duality} has a test, and the test cannot fail. As
\cref{ob:duality} records, \code{Plan.StaleSources} returns the rows the
downward traversal recorded, so \code{TestStaleSourcesAgreesWithPropagation}
compares a quantity with itself. That is a reasonable engineering choice (an
independent upward walk would be a second implementation of the same rule, with
its own bugs) but it means the duality is asserted rather than checked here,
and it is the reason \cref{ob:duality} went unnoticed. The randomised
cross-check \spec{13.7b} recommends is the thing that would catch it, and it
needs a workspace with a prerelease train to do so.

\Cref{thm:g6} is the row whose test predates the normative CCME 2.0.0 fresh-window rule. \code{TestPrereleaseTrainConvergesAfterRelease}
opens with the words ``the reported bug'': a package that entered a beta train
and published \code{1.0.0-beta.0} re-appeared in every subsequent plan as
\code{beta.1}, \code{beta.2}, and so on. That is exactly \cref{ob:pre}, found
in operation rather than on paper, and fixed in the engine. CCME 2.0.0 now
specifies $W_{\rm fresh}$, so a second implementation following the current
text discharges the prerelease work as the engine does.

\Cref{prop:bucket} has vectors and no test because the reference
implementation does not perform the optimisation the vectors are about. Its
traversal is per unit (\code{propagate.go}), which is the literal reading and
is trivially correct. The vectors of Appendix B.7a exist for implementations
that do optimise, and \cref{ob:bucket} adds one more.

\subsection*{Randomised differential validation}

After this report was written, its model was validated differentially
against the reference implementation, most recently in the historical campaign
with the released \code{v1.7.2} binary (the \code{v1.4.0}, \code{v1.6.0}, \code{v1.7.0} and
\code{v1.7.2} campaigns agree per-seed on all 3{,}000 instances, so the figures below
hold for those releases): a $\sim$300-line oracle implementing the report model as it stood against the 1.0.0 snapshot, a seeded generator
of instances (5--20 packages, layered graphs, 4--12 units, mixed tag
states), and a realiser turning each instance into a real repository
for the binary to plan.

On 1{,}500 instances inside the historical premises then called (H1) and (H2), across two disjoint
seed ranges, the two plans agreed
\emph{exactly}: package sets, versions, bumps and channels. On
1{,}000 instances over the full directive grammar the 1{,}097
divergent package-rows classified \emph{completely} into four
predicted classes and none outside them: 240 prerelease-train
discharges and 58 suppression-lapse and downstream-cascade variants
are \cref{ob:pre,ob:g5}
observed at scale; 409 channel-entry patches are \SPEC{}~\S11.4's
deliberate one-patch rule (\diag{W204}), which neither this report's
model nor the companion article's transcribes; and 390 bump-axis
liveness cases arise because the oracle implements the companion
article's compact admission, which drops this report's clause
\ref{a:live}, a transcription gap of the article, caught by the
differential run exactly as intended. The obligations plus these two
recorded omissions form a complete divergence taxonomy for the grammar
exercised, which is the strongest empirical statement this historical campaign
makes about its scope. These records do not validate changed CCME 2.0.0
admission.

The fresh 7 September 2026 E1 campaign used the released
\code{dispat 1.8.2 (linux\_arm64)} binary on 3{,}000 generated instances. The
runner validates Git fixtures and CLI exits. All 1{,}500 core instances agreed
with no invalid sample. The full profile intentionally generates regressive
prereleases: 997 plans were accepted, of which 546 matched the historical
oracle, while 503 complete diagnostic plans were semantically refused with
\diag{E195}. Four refused dossiers had no row difference, so 550 comparisons
were diagnostic-clean. All proposed regression package sets in the 503
dossiers were independently checked, including held proposals. The full
profile had 1{,}914 divergent package rows, assigned by the campaign's heuristic
classifier to PG1 862, PG2 375, PG2x 14, and PG3? 663, with zero unclassified
rows. The canonical manifest binds all four shards to extractor SHA-256
\code{3d5eef3e...11e52}; its paths and hashes identify the four records used
for these totals. Because the oracle
implements the historical model and the labels---especially PG3?---are
heuristic, these counts are taxonomy evidence, not CCME 2.0.0 conformance.

E1 uses tmpfs and pre-run \code{git fsck}; the revised exit checker separates
validated E195 refusals from execution failures. The first stopped attempt
and the 1.8.1 records remain preserved. Seed 580's earlier invalid fixture
remains unresolved; its fresh success is not a product-fix claim, since
planner source is unchanged and the fixture environment changed.

A separate check in the 7 September 2026 campaign ran vectors 133--138 against the released
Dispat 1.8.2 binary; all six passed. The retained raw before/after plan streams
for 134 and 137 confirm package counts, versions and catch-up behaviour. Vector
135 exercises the observable downward eligibility result; it is not an
independent implementation of the upward audit, so the proof remains the
support for full duality. This is selected conformance evidence, not an
exhaustive proof of the specification.


The official image is
\code{yohimik/dispat-alpine@sha256:4c0ad4b8...e7bfd}. The provenance bundle
fixes peeled release tag commit \code{a7011cbcb9e8} and verifies the official
binary against its release asset. All 753 files in the final checksum
manifest and all eight completion/identity markers verified. The same bundle records the other
fresh checks without converting
them into proof claims. E2 repeated orphan recovery twice: no refused tag, the
same recovered version, then an empty third plan. E3 reproduced the established
historical replay profile over five releases per repository (36/42 exact for
\code{google-cloud-node}, 7/64 for \code{conventional-changelog}); it does not
claim those releases followed CCME. E4's two recovery repetitions converged in
24/24 episodes within four runs (31 and 32 total runs) with no structural G3/G5 errors. Its eight
fresh timing cells span 100--10{,}000 packages: medians range from 0.017 s for
the 100-package layered case to 19.354 s for the 10{,}000-package chain, with
15--118 MiB observed maximum RSS. The recorded fixtures passed \code{git fsck};
the pad generator was used only after a small fixture matched the ordinary
generator's commit IDs, refs, messages, and log output. Measurements used a
contended Apple-silicon Docker host capped at two CPUs and 4 GiB, so they are
observations rather than clean comparisons. E5 freshly reran its orphan,
clean concurrent-push and conflicting concurrent-push Dispat cells; their
8, 12, and 14 checks passed. Its nine competitor cells remain historical evidence.
